\chapter{56}



«Hi, Ciro. Alles gut bei dir?»

Trotz fortgeschrittener Stunde klang er frisch.

«Joh, hallo, ja, bestens. Mazi ist ein guter Lehrer.»

«Das kann ich mir vorstellen,» er lachte, «gut, hör mal, bei uns hat sich einiges verändert. David ist bereits zurück. Wir kommen schon am Freitag rauf und werden Besuch mit hoch bringen. Kannst du
bitte das dritte Bett unten bereitstellen und frisch beziehen? Wo die
Wäsche ist, hat dir Manuela gezeigt, nicht?»

«Ja, klar, mache ich.»

«Noch was. Bitte schau, dass die Hütte sauber ist. Die beiden Damen
lieben es ordentlich,» Joh hörte sich überdreht an, «Manuela beauftragte
mich, es dir mitzuteilen. Sie denkt wohl an Mazis gelegentliches Chaos,
hahaha.»

Louis hörte im Hintergrund knapp Manuelas Stimme.

«Na hör mal, was erzählst du da.»

Louis setzte sich auf. Die beiden schafften es innerhalb von Sekunden,
ihn aufzuheitern. Er spielte vergnügt mit.

«Zu Befehl, Chef. Ich werde die Bude auf Hochglanz polieren.»

Jetzt drang zweistimmiges Gelächter laut aus dem Handy. Louis löste das
Gerät vom Ohr und lachte mit.

«Also, Ciro,» sie schienen sich gefasst zu haben, «wir kommen am
Vormittag, ich denk zwischen elf und zwölf. Schlaf gut. Und wenn du uns
brauchst, einfach anrufen, ja?»

«Okay, danke und gute Nacht.»

Louis stellte das Handy aus und legte es neben dem Bett auf den
Bretterboden. Er hatte das Gefühl, bereits ein bisschen dazuzugehören.
Bald würde er David kennenlernen und wer die zweite «Dame» war,
würde sich zeigen. Er drehte sich schwungvoll auf den Bauch. Bevor er
den nächsten Schlafversuch startete, sah er erneut Mazi vor sich.
Er erinnerte sich an ihre kurzen Gespräche zwischen den Schafbesuchen und
an die vielen Fragen, die Mazi in ihm auslösten. An Mazis Aussagen über
seine Eltern. Über seine Kindheit, die ihn zum Strahlen brachte. Von der
er erzählte, als wäre er ein Glückskind der seltenen Art.
«Als Kind kennst du nur, was du kennst,» hatte er schliesslich gesagt,
«meine Eltern hatten es gut zusammen. Erst als Jugendlicher verstand ich
nach und nach, dass in andern Familien so einiges los ist. Manchmal
empfand ich meine Familie nahezu als langweilig. Ich konnte meinen
Freunden keine Katastrophengeschichten erzählen, beim Verfluchen der
Eltern nicht mithalten,» dann hatte Mazi den Kopf geschüttelt und die
Hände verworfen, «dadurch lernte ich zuzuhören, ganz selten erfand ich
irgend etwas Halbwildes, um nicht abzufallen. Oder schmückte meine
Afrikabesuche aus. Das interessierte zwar die Jungs mässig, die Mädchen
aber waren begeistert und wären gerne zu meinen Verwandten mitgekommen,»
dann hatte er wieder sein weisses Blitzlachen gezeigt, «heute find ich’s
lustig und bin einfach nur dankbar für meine tollen Eltern.»
Wieder waren die Worte ohne grosses Aufhebens aus ihm heraus gepurzelt. 
Mazis Geschichte lebte von unaufgeregter Einfachheit. Einmal mehr
wunderte sich Louis über die Offenheit eines ihm weitgehend Fremden. Er
war in eine durch und durch wundersame Gesellschaft geraten.
Mit dem Gedanken schloss Louis die Augen. Er schlief augenblicklich ein.
